\documentclass[12pt,a4paper]{article}

\usepackage[T1]{fontenc}
\usepackage[utf8]{inputenc}
\usepackage[romanian]{babel}
\usepackage{geometry}

\geometry{margin=2.5cm}

\title{Tutoriat 3\\Cautare Adversariala - Recapitulare si Exercitii}
\author{}
\date{}

\begin{document}

\maketitle

\section{Recapitulare}

În acest tutoriat au fost prezentate principalele concepte utilizate în căutarea adversarială.

\subsection{Concepte fundamentale}

Am studiat:

\begin{itemize}
\item agenții inteligenți;
\item mediile de lucru;
\item jocurile cu doi jucători;
\item jocurile cu sumă zero;
\item formularea unei probleme de joc;
\item arborele de joc.
\end{itemize}

\subsection{MiniMax}

Algoritmul MiniMax:

\begin{itemize}
\item presupune că ambii jucători joacă optim;
\item utilizează noduri MAX și MIN;
\item propagă valorile de la frunze către rădăcină;
\item determină mutarea optimă.
\end{itemize}

Reguli:

\[
MAX \rightarrow \max
\]

\[
MIN \rightarrow \min
\]

\subsection{Limitările MiniMax}

Principala problemă este creșterea exponențială a numărului de stări.

\[
O(b^d)
\]

unde:

\begin{itemize}
\item $b$ este factorul de ramificare;
\item $d$ este adâncimea arborelui.
\end{itemize}

\subsection{Funcții de evaluare}

Atunci când arborele nu poate fi explorat complet:

\begin{itemize}
\item se limitează adâncimea de căutare;
\item se utilizează funcții de evaluare;
\item stările sunt estimate în loc să fie evaluate exact.
\end{itemize}

\subsection{Alpha-Beta Pruning}

Alpha-Beta este o optimizare a algoritmului MiniMax.

Avantaje:

\begin{itemize}
\item elimină ramurile inutile;
\item reduce timpul de execuție;
\item păstrează aceeași soluție optimă;
\item permite explorarea unor arbori mai adânci.
\end{itemize}

Condiția de retezare:

\[
\alpha \geq \beta
\]

\subsection{Idei de reținut}

\begin{enumerate}
\item Jocurile sunt modelate prin arbori de joc.
\item MAX încearcă să maximizeze scorul.
\item MIN încearcă să minimizeze scorul.
\item MiniMax găsește strategia optimă.
\item Funcțiile de evaluare permit limitarea adâncimii.
\item Alpha-Beta Pruning accelerează MiniMax.
\item Alpha-Beta produce aceeași soluție ca MiniMax, analizând mai puține noduri.
\end{enumerate}












\section{Exerciții}

\subsection{Exercițiul 1}

Considerăm arborele:

\begin{verbatim}
          MAX
         /   \
       MIN   MIN
      / \    / \
     3   5  2   9
\end{verbatim}

Determinați valoarea fiecărui nod și valoarea rădăcinii folosind algoritmul MiniMax.

\subsection{Exercițiul 2}

Considerăm arborele:

\begin{verbatim}
              MAX
           /       \
         MIN       MIN
        /  \       /  \
       8    4     6    2
\end{verbatim}

Determinați:

\begin{itemize}
\item valorile nodurilor MIN;
\item valoarea rădăcinii;
\item mutarea optimă pentru MAX.
\end{itemize}

\subsection{Exercițiul 3}

Pentru arborele:

\begin{verbatim}
                MAX
             /       \
           MIN       MIN
          /  \       /  \
        MAX  MAX   MAX  MAX
        / \  / \   / \  / \
       1  4 2  7  3  5 6  8
\end{verbatim}

Aplicați pas cu pas algoritmul MiniMax.

\subsection{Exercițiul 4}

Considerăm un joc cu următoarea funcție de utilitate:

\begin{itemize}
\item victorie MAX = +1;
\item remiză = 0;
\item victorie MIN = -1.
\end{itemize}

Ce valoare va avea o stare terminală în care MAX pierde?

Dar una în care partida se termină la egalitate?

\subsection{Exercițiul 5}

Fie:

\[
b = 5 (branching factor)
\]

și

\[
d = 4 (depth)
\]

Determinați numărul aproximativ de noduri explorate de algoritmul MiniMax.

\subsection{Exercițiul 6}

Explicați diferența dintre:

\begin{itemize}
\item joc determinist;
\item joc stochastic.
\end{itemize}

Dați câte un exemplu pentru fiecare situație.

\subsection{Exercițiul 7}

Explicați diferența dintre:

\begin{itemize}
\item joc cu informație completă;
\item joc cu informație incompletă.
\end{itemize}

Dați exemple de jocuri pentru fiecare categorie.

\subsection{Exercițiul 8}

Considerăm arborele:

\begin{verbatim}
          MAX
         /   \
       MIN   MIN
      / \    / \
     3   7  2   8
\end{verbatim}

Indicați ce ramuri pot fi eliminate utilizând Alpha-Beta Pruning.

Justificați răspunsul.

\subsection{Exercițiul 9}

Care este rolul parametrilor:

\begin{itemize}
\item alpha;
\item beta.
\end{itemize}

În ce condiție se realizează retezarea unei ramuri?

\subsection{Exercițiul 10}

Explicați avantajele utilizării Alpha-Beta Pruning față de MiniMax.

\subsection{Exercițiu recapitulativ}

Completați spațiile libere:

\begin{enumerate}
\item MiniMax este utilizat pentru probleme de .................... .
\item MAX încearcă să .................... valoarea unei stări.
\item MIN încearcă să .................... valoarea unei stări.
\item Jocurile studiate sunt de regulă jocuri  .................... .
\item Funcțiile de evaluare sunt folosite atunci când .................... .
\item Alpha-Beta Pruning produce aceeași soluție ca .................... .
\item O ramură poate fi eliminată atunci când .................... .
\end{enumerate}


\end{document}